% Source: Supplement: shortest path formulation of difference constraints.
% Review and validation: docs/validation.md.
\subsection{差分约束}
\label{algo:difference-constraints}
\AlgorithmInfo
    {判定形如 $x_u-x_v\leq w$ 的不等式组是否有解，并构造一组整数解。}
    {本实现使用 SPFA，最坏时间 $O(nm)$，空间 $O(n+m)$。}
    {约束常数为整数，绝对值小于 $4\times10^{18}$；可行势值绝对值也须小于该界。}

\subsubsection*{原理与连边方向}
把不等式改写为 $x_u\leq x_v+w$，恰好对应最短路的三角不等式，所以从 $v$ 向 $u$ 连权为 $w$ 的边。所有点从超级源可达时，若图有负环，把环上的不等式相加会得到 $0<0$ 的矛盾；没有负环时，最短距离本身就是一组满足所有约束的赋值。

本模板将所有初始距离设为 $0$，等价于补上超级源的零权边，确保不漏掉不连通部分的矛盾。只从某一个变量出发检查负环是不够的。

常见变形：$x_u-x_v\geq w$ 改写为 $x_v-x_u\leq-w$；等式拆成两个方向的不等式。整数变量的严格小于 $x_u-x_v<w$ 可转为 $\leq w-1$，实数变量不能直接减一。

\subsubsection*{使用教程}
先复制 SPFA。\InlineCode{le(u,v,w)} 表示 $x_u-x_v\leq w$；\InlineCode{ge} 表示大于等于，\InlineCode{eq} 表示等于。\InlineCode{solve()} 返回真后读取 \InlineCode{val[u]}，为假表示约束矛盾。加约束后可重新求解，多测调用 \InlineCode{init(n)}。
\begin{lstlisting}[style=mtycpp,numbers=none]
DifferenceConstraints dc;
dc.init(3);
dc.le(1,2,3);
dc.ge(2,3,1);
dc.eq(1,3,2);
bool ok=dc.solve();
if(ok){long long x=dc.val[1];}
\end{lstlisting}

\needspace{7\baselineskip}
\subsubsection*{模板代码}
\CodeFile{code/graphs/difference_constraints.cpp}

若仅包含差值约束，所有变量同时平移不影响可行性，可按需减去最小值使解非负。该接口求的是一组可行解，不保证某个附加目标最优；若要求最大差值上界、固定变量或最小总和，应另行推导目标与源点的关系。
